\section*{கணங்களின் அளவு பற்றிய பதிப்புக் குறிப்புகள்}
இந்தக் குறிப்புகள் மூல நூலின் பகுதிகள் அல்ல.

\textbf{TA-SIZ-LOCAL-001.} ஜோடியாக்கச் சார்புகள் பற்றிய பிரிவில்,
$k$ ஆவது முக்கோண எண் “$k$ ஐ விடச் சிறிய இயல் எண்களின் கூட்டுத்தொகை”
என்றும், அதன் வாய்பாடு $k(k+1)/2$ என்றும் உறையவைக்கப்பட்ட ஆங்கில மூலம்
கூறுகிறது. இந்த நூலில் இயல் எண்கள் பூச்சியத்திலிருந்து தொடங்குவதால்,
$k$ ஐ விடச் சிறிய இயல் எண்களின் கூட்டுத்தொகை உண்மையில்
$k(k-1)/2$ ஆகும். ஆனால் அந்தப் பிரிவின் $0,1,3,6,\ldots$ என்ற
வரிசையும் அதன் ஜோடியாக்க வாய்பாடும் $k(k+1)/2$ என்ற
வழக்கமான பூச்சியம்-அடிப்படை முக்கோண எண் குறியிடலைப் பயன்படுத்துகின்றன.
மூல வாக்கியமும் வாய்பாடுகளும் உடல் உரையில் வைக்கப்பட்டுள்ளன;
இது குறியெண் வாய்பாட்டின் மதிப்புகளை மாற்றுவதில்லை.

\textbf{TA-SIZ-LOCAL-002.} மாற்று ஜோடியாக்கத்தின் முதலாவது படியில்,
“இரண்டாவது” ஜோடியை உறையவைக்கப்பட்ட ஆங்கில மூலம்
$\tuple{0,2}$ என்கிறது; உடனுள்ள அட்டவணையில் அது
$\tuple{0,1}$ ஆக உள்ளது. அட்டவணைக்கும் தொடர்வரிசைக்கும் பொருந்துமாறு
தமிழ் உடல் $\tuple{0,1}$ எனத் திருத்துகிறது.

\textbf{OLSIZ-001.} முழுக்களைப் பட்டியலாக்கும் அட்டவணையின் தலைப்பு
$f(7)$ வரை செல்கிறது; அதற்கு ஒத்த வாய்பாடு
$-\lceil 6/2\rceil=-3$. ஆனால் உறையவைக்கப்பட்ட ஆங்கில மூலத்தின்
மதிப்பு வரிசை $3$ உடன் நின்று $-3$ ஐ விடுகிறது.
தமிழ் அட்டவணை $f(7)$ இன் மதிப்பாக $-3$ ஐச் சேர்க்கிறது.

\textbf{OLSIZ-002.} இணைமுடிவுறு உட்கணத்தின் வரையறையில்,
உறையவைக்கப்பட்ட ஆங்கில மூலத்திலுள்ள “ஒரு முடிவுறு கணம் $\Nat$ இன்
நிரப்பி” என்ற சொற்றொடர், முடிவுறு கணத்திற்கும் $\Nat$ க்கும் உள்ள
உறவை விடுகிறது. உடனுள்ள $\Nat\setminus A$ முடிவுறும் என்ற
சமானநிபந்தனைக்கு ஏற்ப, தமிழ் உடல் இதனை $\Nat$ இன் ஒரு முடிவுறு
உட்கணத்திற்கு, இயல் எண்களின் கணத்திற்குள்ளான நிரப்பி எனத் தெளிவாக்குகிறது.

\textbf{OLSIZ-003.} மாற்று ஜோடியாக்கக் கட்டுமானத்தின் அடுத்த படியில்
$\tuple{2,m}$ என்பதை ஆங்கில மூலம் இருமுறை அடுத்தடுத்து அச்சிடுகிறது.
பின்னுள்ள நிறைவடைந்த பட்டியலில் கிடைவரிசைகள் $2$, $3$ என்று
தொடர்வதால், தமிழ் உடல் இரண்டாவதை $\tuple{3,m}$ எனத் திருத்துகிறது.

\textbf{OLSIZ-004.} முதன்மைக் குறைத்தல் நிறுவலில்,
$f(Z)$ இன் வெளியீட்டை $s_k$ என்று உறையவைக்கப்பட்ட ஆங்கில மூலம்
பெயரிடுகிறது; ஆனால் $k$ அங்கு கட்டுப்படுத்தப்படவோ தீர்மானிக்கப்படவோ
இல்லை. அடுத்த உறுப்புவாரி நிபந்தனை $Z$ இன் இயல்புக்குறித் தொடரைத்
தனித்த முறையில் தீர்மானிக்கிறது. தமிழ் உடல் அந்த வெளியீட்டையும் அதன்
உறுப்புகளையும் ஒரே கட்டுப்பட்ட பெயரான $s$, $s(n)$ கொண்டு குறிக்கிறது.

\textbf{OLSIZ-005.} அதே பிரிவின் எதிரெடுத்துக்காட்டில்,
$h\colon\PosInt\to\Bin^\omega$ என அறிவித்த பிறகு, $h(n)$ ஐ
$n$ பூச்சியங்களைக் கொண்ட முடிவுறு தொடராக உறையவைக்கப்பட்ட ஆங்கில மூலம்
காட்டுகிறது. அது $\Bin^\omega$ இன் உறுப்பு அல்ல. தமிழ் உடல்,
$n$ பூச்சியங்களுக்குப் பின் முடிவுறாமல் ஒன்றுகள் வரும் தொடரைப்
பயன்படுத்துகிறது; ஆகவே ஒவ்வொரு வெளியீடும் உண்மையாகவே
$\Bin^\omega$ இல் உள்ளது.

\textbf{OLSIZ-006.} எண்ணொத்த தன்மை பற்றிய நிறுவலில்,
$f\colon A\to B$ என்ற இருபுறச் சார்பை முன்கொண்ட பிறகு,
$A=\emptyset$ எனில் $B=\emptyset$ என்பதை விளக்கும் இரண்டு
நிபந்தனைக் கிளைகளிலும் உறையவைக்கப்பட்ட ஆங்கில மூலம்
$g(x)=y$ என்கிறது. அந்த நிலையில் தேவையான சார்பு $f$ ஆகும்;
$g$ என்பது பின்னர் பயன்படுத்தப்படும் பட்டியலாக்கம்.
தமிழ் உடல் இரண்டு இடங்களிலும் $f(x)=y$ எனத் திருத்துகிறது.

\textbf{OLSIZ-007.} காண்டோரின் தேற்றத்தின் விரிவான நிறுவலில்,
ஏதேனும் ஒரு $x\in A$ க்கு $g(x)\ne\overline A$ எனக்
காட்ட வேண்டிய இடத்தில், உறையவைக்கப்பட்ட ஆங்கில மூலம்
“ஒவ்வொரு $x\in\overline A$ க்கும்” என்கிறது.
முந்தைய ஏதேனும்-ஒன்றான தேர்வுக்கும் அடுத்த முடிவுக்கும் பொருந்துமாறு,
தமிழ் உடல் “ஒவ்வொரு $x\in A$ க்கும்” எனத் திருத்துகிறது.

\textbf{OLSIZ-008.} மாற்று எண்ணிடத்தகாமைப் பிரிவின் மூலைவிட்ட
அட்டவணையை அறிமுகப்படுத்தும்போது, $s_n(m)$ என்பதை உறையவைக்கப்பட்ட
ஆங்கில மூலம் “$m$ ஆவது தொடரின் $n$ ஆவது இலக்கம்” என்கிறது.
ஆனால் உடனுள்ள அட்டவணையில் $s_n$ என்பது $n$ ஆவது கிடைவரிசைத் தொடரும்,
$m$ என்பது அதன் நெடுவரிசை இலக்கமும் ஆகும்.
தமிழ் உடல் இதனை “$n$ ஆவது தொடரின் $m$ ஆவது இலக்கம்” எனத் திருத்துகிறது.

\textbf{OLSIZ-009.} அதே நிறுவலின் மூலைவிட்டக் கட்டுமானத்தில்,
ஒவ்வொரு $1$ ஐயும் $0$ ஆகவும் ஒவ்வொரு $1$ ஐயும் $0$ ஆகவும்
மாற்றுவதாக ஆங்கில மூலம் இருமுறை ஒரே மாற்றத்தைச் சொல்கிறது.
அடுத்த காட்சிவாய்பாடு காட்டும் முழுமாற்றத்திற்கு ஏற்ப,
தமிழ் உடல் ஒவ்வொரு $1$ ஐயும் $0$ ஆகவும் ஒவ்வொரு $0$ ஐயும்
$1$ ஆகவும் மாற்றுகிறது.

\textbf{OLSIZ-010.} மாற்றுக் குறைத்தல் நிறுவலில்,
$f(N)$ இன் வெளியீட்டை மீண்டும் $s_k$ என்று உறையவைக்கப்பட்ட ஆங்கில
மூலம் பெயரிடுகிறது; இங்கும் $k$ கட்டுப்படுத்தப்படவில்லை.
உறுப்புவாரி நிபந்தனை $N$ இன் இயல்புக்குறித் தொடரைத் தனித்த முறையில்
தீர்மானிப்பதால், தமிழ் உடல் வெளியீட்டையும் அதன் உறுப்புகளையும்
$s$, $s(n)$ என்ற ஒரே கட்டுப்பட்ட பெயரால் குறிக்கிறது.

\textbf{TA-SIZ-LOCAL-003.} முதன்மை மற்றும் மாற்றுக் குறைத்தல்
பிரிவுகளில் உள்ள இரு வேறுபட்ட பயிற்சிகளுக்கும் உறையவைக்கப்பட்ட ஆங்கில
மூலம் ஒரே முழு அடையாளமான
\texttt{sfr:siz:red:prob:nat-nat} ஐ வழங்குகிறது. இரு பிரிவுகளையும்
ஒரே நூலில் சேர்க்கும்போது இது இரட்டைச் சுட்டியாகிறது.
மாற்றுப் பிரிவின் பெயரிடைவெளிக்கு ஏற்ப, தமிழ் மாற்றுப் பயிற்சியின்
அடையாளத்தை \texttt{sfr:siz:red-alt:prob:nat-nat} எனத் திருத்துகிறது;
முதன்மைப் பயிற்சியின் அடையாளம் மாறவில்லை.
